% !TEX program = lualatex
\documentclass[a4paper,12pt]{article}
\usepackage{fontspec}
\usepackage{polyglossia}
\usepackage{amsmath}
\usepackage{microtype}
\usepackage{unicode-math}
\usepackage[left=15mm,right=15mm,top=15mm,bottom=20mm]{geometry}
\usepackage{graphicx}
\usepackage{caption}
% \usepackage{tabularx}
\usepackage{wrapfig}
\usepackage{mdframed}
\usepackage{setspace}
\usepackage[version=4]{mhchem}


\setmainlanguage{russian}
\setotherlanguage{english}
% Переключение: \textenglish{...} или \begin{english}...\end{english}
\setmainfont{Times New Roman}
\setmonofont{Source Code Pro}
% \setmathfont{Libertinus Math}
% \setmathfont{STIX Two Math}
\setmathfont{Cambria Math}
\setlength{\parindent}{10mm}
\setlength{\parskip}{2mm}

\pretocmd{\subsection}{\clearpage}{}{}

\begin{document}

		\paragraph*{Электрический заряд.} \

        Тела, которые действуют на окружающие предметы электрическими силами,
		являются наэлектризованными или заряженными; на этих телах находятся
		заряды.

		Посредством трения можно зарядить любое тело.

		Электрические заряды могут переходить с одного тела на другое.

		Прикасаясь заряженным телом к незаряженным, им передается электрический
		заряд.

		Электрические заряды могут также перемещаться по телу.

		Вещества, по которым электрические заряды легко перемещаются, называют
		\textit{проводниками}, не обладающие этим свойством, называют
		\textit{диэлектриками (изоляторами)}.

		\textit{Разделение на проводники и диэлектрики условно.}

		Большинство веществ природы принадлежит к группе \textit{полупроводников}
		(по своим свойствам занимают промежуточное положение между очень хорошими
		проводниками и очень хорошими диэлектриками).

		Изолирующие свойства вещества зависят также от его состояния и могут
		сильно изменяться.

		Когда в каком-либо теле происходит перемещение зарядов, говорят, что в
		этом теле имеется электрический ток.

		



\end{document}